% Generated draft: requires mathematical and source review.
\documentclass[11pt]{article}
\usepackage[margin=25mm]{geometry}
\usepackage{fontspec}
\setmainfont{Latin Modern Roman}
\usepackage{amsmath,amssymb,mathtools,unicode-math}
\setmathfont{Latin Modern Math}
\usepackage{enumitem,needspace,xurl,hyperref}
\setlength{\parindent}{0pt}
\setlength{\parskip}{5pt}
\setlength{\emergencystretch}{3em}
\newcommand{\N}{\mathbb{N}}
\newcommand{\Z}{\mathbb{Z}}
\newcommand{\Q}{\mathbb{Q}}
\newcommand{\R}{\mathbb{R}}
\newcommand{\C}{\mathbb{C}}
\newcommand{\F}{\mathbb{F}}
\newcommand{\A}{\mathbb{A}}
\newcommand{\PP}{\mathbb{P}}
\newcommand{\E}{\mathbb{E}}
\newcommand{\eps}{\varepsilon}
\newcommand{\dd}{\,\mathrm{d}}
\begin{document}
\section*{UnsolvedMath ID: 1100808}
\textbf{Source problem code:} catalogue entry
\subsection{Definitions and mathematical statement}
For a finitely generated group $G$, use a word metric from a finite generating set. Its asymptotic dimension satisfies $\operatorname{asdim}G\leq n$ if for every $R>0$ there is a uniformly bounded cover of $G$ that is the union of $n+1$ families, each consisting of pairwise $R$-disjoint sets. Uniform boundedness means one finite upper bound for the diameters of all sets in that cover; $R$-disjoint means the distance between distinct sets is greater than $R$. The asymptotic dimension is the least such $n\geq0$, or $\infty$ if none exists, and is independent of the finite generating set.

\textbf{Questions.} Determine the asymptotic dimensions of CAT(0) groups, $\operatorname{Out}(F_n)$, mapping class groups, nonuniform lattices, and Thompson's group. Here a CAT(0) group conventionally means a group admitting a proper cocompact isometric action on a proper CAT(0) space, and $\operatorname{Out}(F_n)=\operatorname{Aut}(F_n)/\operatorname{Inn}(F_n)$ for a free group of finite rank $n$. The source does not specify the surface parameters for mapping class groups, the ambient groups for lattices, or which Thompson group in its main sentence.

Does there exist a group of finite type with infinite asymptotic dimension? Two distinct interpretations are relevant: type $F$ means having a finite CW model for $K(G,1)$, while type $F_{\infty}$ means having a $K(G,1)$ with finite skeleton in every dimension. A $K(G,1)$ is a connected CW complex with fundamental group $G$ and contractible universal cover.

The supplied summary reports finite asymptotic dimension for mapping class groups and nonuniform lattices in semisimple Lie groups, and infinite asymptotic dimension for Thompson's group $F$. It reports Thompson's groups $F$ and $V$ as type-$F_{\infty}$ examples answering that interpretation of the existence question affirmatively, while describing the stronger type-$F$ interpretation as still open to the summary author's knowledge. It reports finiteness still unknown for general CAT(0) groups and general $\operatorname{Out}(F_n)$, with the special bound by cubical dimension for cubulated CAT(0) groups. These finiteness statements do not supply exact numerical dimensions for every group in the requested classes.
\subsection{Short English statement}
Determine large-scale dimensions for the listed group families, and ask whether a homotopically finite group can have infinite asymptotic dimension. The supplied progress distinguishes finite-dimensionality from exact values and type $F$ from type $F_{\infty}$.
\subsection{Sources}
\begin{itemize}
\item M. Bestvina, \emph{Questions in Geometric Group Theory} (2004), Question 8.8, PDF page 17: \url{https://www.math.utah.edu/~bestvina/eprints/questions-updated.pdf}; manual transcription, accessed 2026-07-29.
\item Supplied research-summary attributions: Bestvina--Bromberg--Fujiwara (2015), mapping class groups; Dranishnikov--Sapir, Thompson's $F$; Ji (2004) and Osin (2005), nonuniform lattices; Wright (2012), cubical dimension bound; Bestvina--Guirardel--Horbez (2022), boundary amenability of $\operatorname{Out}(F_n)$. No titles or theorem identifiers are supplied.
\item Local registry \texttt{data/unsolvedmath/problems.json}, permanent ID $1100808$, source code catalogue entry, \emph{Questions in Geometric Group Theory --- Q 8.8}. The record labels this partially solved and PARTIAL-PROGRESS, with dated source-status and rights-review qualifications.
\end{itemize}
\end{document}
